\subsection{二つの極限成分を表す一つの係数}
\label{sec:joint-coefficient}
直前の実現で用いた係数選択を詳述する。この節の分解は全て、固定した補助確率 $Q$ の下で、
元の価格 $S$ に対して行う。零始点の正準分解を $S-S_0=M^S+A^S$ とし、
実現済み近似積分の係数と成分を
\[
 M_n=H_n\cdot M^S,\qquad A_n=H_n\cdot A^S
\]
と書く。前述の成分級数の構成により、正則な極限 $M,A$、$M_n$ の局所一様収束、
$A_n$ の変動収束、および各有限 $T$ について概確実な
\[
 \sum_n\Var_{[0,T]}(A_{n+1}-A_n)<\infty
\]
が得られている。さらに共通の joint 停止列が得られ、各座標で停止したM成分の終端は、
近似列の添字について一様に $L^2(Q)$ で有界である。

\begin{lemma}[二つの制御測度に対する同時係数選択]\label{lem:joint-coefficient}
この条件の下で、一つの実数値予測可能係数 $H$ が存在し、各 joint 座標上の完成 M 積分と
FV 積分は、それぞれ $M,A$ の対応する停止と一致する。
energy measure と variation measure の間の絶対連続性は仮定しない。
\end{lemma}
\begin{proof}
まず決定論的な $T>0$ を固定し、元価格の経路的変動測度を
$v_T(\omega)=|d(A^S)^T|(\omega)$ と書く。
\[
 v_T(\omega)([0,\infty))+\Var_{[0,T]}A_0+
 \sum_n\Var_{[0,T]}(A_{n+1}-A_n)
\]
を支配する有限可測確率変数 $R_T\geq0$ を選ぶ。既に示した経路的総和可能性から選べる。
構成では、総和可能な M 成分の包絡も $R_T$ に加えてよく、その場合は同じ重みを全成分の評価に使える。
$a_T=(1+R_T)^{-2}$ と置き、零例外集合では $1$ とする。これは概確実に正である。
予測可能 sigma-field 上の有限測度を
\[
 \nu_T(B)=E_Q\left[a_T\int1_B(t,\omega)\,dv_T(\omega)(t)\right]
\]
で定める。変更するのは FV の制御測度だけであり、M 成分と energy の評価は $Q$ の下に保つ。

各近似積分の経路的密度等式より、
\[
 \begin{aligned}
 \int|H_0|\,d\nu_T&=E_Q[a_T\Var_{[0,T]}A_0]<\infty,\\
 \sum_n\int|H_{n+1}-H_n|\,d\nu_T
 &=E_Q\left[a_T\sum_n\Var_{[0,T]}(A_{n+1}-A_n)\right]<\infty.
 \end{aligned}
\]
実際、$d(A_{n+1}-A_n)=(H_{n+1}-H_n)dA^S$ であり、密度 $f$ を持つ符号付き測度の
全変動は、元の全変動測度に密度 $|f|$ を掛けたものである。
Tonelli と $L^1$ の完備性から、係数の列全体は $h_T$ に
$\nu_T$-a.e. かつ $L^1(\nu_T)$ で収束する。
経路ごとにも、最初の項は $v_T(\omega)$ 可積分で、隣接項の $L^1$ 距離の和は有限である。
従って一つの $Q$-零集合の外で、$v_T(\omega)$-a.e. および
経路的な $L^1(v_T(\omega))$ 収束が成り立つ。

可算個のjoint終端座標について、$\sup_n\|M_n(\rho_r)\|_2\leq b_r$ となる
$b_r>0$ を取る。ベクトル $(2^{-(r+1)/2}M_n(\rho_r)/b_r)_r$ は終端 $L^2(Q)$ 空間の
Hilbert直和で有界である。弱部分列コンパクト性とMazurの補題により、
この直和で強収束するforward凸平均を選べる。部分列の重みを元の添字に戻しても台は前進するので、
一つの重み族 $w_{n,k}$ によって全ての終端座標がCauchyとなる。
$\bar H_n=\sum_kw_{n,k}H_k$ と置き、$\bar M_n$ にも同じ重みを使う。
$r$ 番目の座標のpredictable energy measureを $q_r$ とすると、完成積分の等長性より
\[
 \|\bar H_n-\bar H_m\|_{L^2(q_r)}
 =\|\bar M_n(\rho_r)-\bar M_m(\rho_r)\|_{L^2(Q)}
\]
である。ここで $\rho_r$ は決定論的 horizon も含めた有限な座標停止時刻であり、
$\bar M_n$ は同じ重みの M 成分である。
従って係数は可算個の全ての $L^2(q_r)$ で Cauchy である。
一つの狭義増加な $f$ を
\[
 \|\bar H_{f(j+1)}-\bar H_{f(j)}\|_{L^2(q_r)}\leq2^{-j}
 \quad(r\leq j)
\]
となるように選ぶ。$q_r$ は有限なので、Cauchy--Schwarz と Tonelli により
$q_r$-a.e. に点ごとの収束が成り立つ。この列が $\mathbb R$ 内で収束する点ではその極限を、
それ以外では零を $H$ と定める。収束集合は予測可能なので $H$ も予測可能である。
$L^2$ の完備性より、各 $r$ で $\bar H_{f(j)}\to H$ は $L^2(q_r)$ 収束でもある。

この同じ定義は、$q_r$ が零となる部分も含めて FV 側の係数を決める。
$H_n\to h_T$ となる点では、forward 凸平均の性質から
\[
 |\bar H_n-h_T|\leq\sup_{k\geq n}|H_k-h_T|\longrightarrow0
\]
である。従って $H=h_T$ が $\nu_T$-a.e. に成立し、ほとんど全ての経路で
$v_T(\omega)$-a.e. にも成立する。特に $H\in L^1(\nu_T)$ であり、
経路ごとに $v_T$ に関して可積分である。
joint source prefix は決定論的 horizon より前での追加停止なので、その変動測度は $v_T$ の制限である。
同じ可積分性がその座標でも成立する。これで一つの係数の $L^2(q_r)$ と $L^1$ の所属がそろう。

完成 M 積分の等長性と Doob 不等式により、$\bar H_{f(j)}$ の M 積分は、
座標 $r$ 上で $H$ の M 積分へ収束する。一方、forward 凸平均は既存の極限 $M$ を保つ。
確率収束の極限の一意性から、この積分は $M^{\rho_r}$ に一致する。
FV 側では、ほとんど全ての経路で
\[
 \Var_{[0,T]}\bigl(H_n\cdot A^S-H\cdot A^S\bigr)
 =\int|H_n-H|\,dv_T\longrightarrow0
\]
である。同じ積分列は変動距離で $A^T$ に収束するので、零始点性と極限の一意性から
$H\cdot(A^S)^T=A^T$ を得る。
この符号付き測度の等式を $(0,\rho_r]$ に制限すれば、停止時の jump も含めた
joint 座標での FV 同定を得る。可算個の座標と右連続性により、過程等式の零集合を一つにできる。
従って同じ $H$ が両成分を実現する。
\end{proof}
異なる外側の停止段階で選んだ係数同士の一致は要求しない。
\ref{sec:integral-calculus}節では、元測度へ戻した後、互いに素な停止区間上でそれらを貼り合わせる。
補題\ref{lem:joint-coefficient}はその入力となる停止利得の実現を与えるものであり、
係数選択の段階で積分領域の閉性を仮定していない。
